% =====================================================================
% Section V edits carrying the hardware runs of HW_RUNS.md into the draft.
% Each block says where it goes and what it replaces.  Numbers are from
% hw_runs/; SECV_HW_TEXT.md says which measurement backs which sentence.
% =====================================================================


% ---------------------------------------------------------------------
% [1] V-A, Digital twin testbed.  Replaces the sentence ending
%     "...reducing the expected sampling cost by a factor of approximately
%     nine for k=1."   (P_z(k)/mu_z is 5.9 as measured; 9 is the mu->0 bound.)
% ---------------------------------------------------------------------
For small $\mu_z$, \eqref{eq:success} gives $P_z(k)\simeq(2k+1)^2\mu_z$, so
$(2k+1)^2=9$ bounds the reduction in sampling cost at $k=1$; the ratio
$P_z(k)/\mu_z$ measured over the zones of this evaluation is $5.9$.


% ---------------------------------------------------------------------
% [2] V-B.  New paragraph, placed after the simulated-noise discussion and
%     before the routing paragraph.  This is the passage that carries the
%     hardware claim of the abstract, Sec. I and the conclusion.
% ---------------------------------------------------------------------
The same construction is executed on quantum processors: the superconducting
devices ibm\_boston and ibm\_kingston, and the trapped-ion device IonQ Forte
Enterprise, whose all-to-all connectivity needs no routing and executes the
inter-cell circuit in $196$ two-qubit gates against $787$ after routing onto a
heavy-hex lattice. Every reported assignment satisfies the original
constraints on every device and in every run, since the acceptance test of
Section~\ref{sec:distribution} is applied to the decoded assignment rather
than to the internal flags. Without amplification, ibm\_boston reproduces the
accepted law of \eqref{eq:conditional-exponential} across six demand settings:
over $10{,}000$ shots the distance to the exact law is $0.024$ to $0.049$,
acceptance lands within $0.92$ to $1.08$ of its exact value, and an accepted
sample carries $0.83$ to $0.90$ of $J^\star$ against an exact $0.84$ to $0.91$.
The utility rotations and the decoded acceptance test therefore operate as
intended on a device.

Amplification is bounded by circuit size rather than by the construction. The
marked branch decays by $0.0022$ to $0.0042$ per routed two-qubit gate over
circuits of $364$ to $1868$ gates, which places the break-even of one round
near $400$ routed gates: a three-UE circuit of $364$ gates raises acceptance
from $0.298$ without amplification to $0.337$ with it, whereas circuits of
$580$ gates and above stay below their own baselines. Selection survives well
past that point. At $1868$ routed gates, half the partition budget in the
counting of \eqref{eq:zone_one_pass_cost}, the surviving signal is $1.5\%$ and
the optimal assignment still holds $0.222$ of the accepted mass against
$0.016$ for a random output. Five physical layouts of that circuit agree
within $0.209$ to $0.251$. Job identifiers, device settings and raw counts are
provided in \cite{scqfrepo}.


% ---------------------------------------------------------------------
% [3] V-B, routing paragraph.  Append to the sentence that says the budget
%     counts logical oracle gates before routing.
% ---------------------------------------------------------------------
The decay measured above places the break-even of one amplification round near
$355$ gates in this counting on a present device, so the budget bounds the
circuit a zone must fit into rather than the size at which a present device
returns an undistorted accepted law; it is met in qubit width, which the zone
circuits of Section~\ref{subsec:eval_network} occupy at $37.7$ to $43.0$
qubits against the $156$ of the devices used here.


% ---------------------------------------------------------------------
% [4] V-B, after the simulated-noise table.  What the noise sweep does and
%     does not cover.
% ---------------------------------------------------------------------
These conditions model gate and readout errors. Measured against the hardware
runs, a noise model carrying the device's own gate errors overestimates the
surviving signal by a factor of $3.4$ at $364$ routed gates and $55$ at
$1868$, because the loss that grows with circuit size is the idling of qubits
through a deep circuit rather than the gates themselves. The sweep therefore
characterizes the effect of gate errors, and the hardware runs the operating
point.


% ---------------------------------------------------------------------
% [5] V-D, after the sentence "Parallel counters reduce the accumulation
%     depth, while the remaining constraint and assignment operations set a
%     floor on the achievable latency."
% ---------------------------------------------------------------------
Depth is also what limits the executable size on present hardware: the
$787$-gate circuit of Section~\ref{subsec:eval_local_law} loses its entire
marked branch, from an acceptance of $0.125$ to $0.034$, when dynamical
decoupling is disabled and its qubits idle through the same depth. Reducing
the accumulation depth is therefore a present-device measure as well as a
latency one.


% ---------------------------------------------------------------------
% [6] Table caption fix, V-B noise table.
%     Was: "...with 2000 shots per condition..."
% ---------------------------------------------------------------------
\caption{Inter-cell circuit under simulated noise, with 10\,000 shots per
condition and one amplification round. TVD is computed after the decoded
feasibility check.}


% ---------------------------------------------------------------------
% [7] Optional table for V-B, if the hardware runs are given their own float
%     rather than only prose.
% ---------------------------------------------------------------------
\begin{table}[!t]
\centering
\caption{Hardware execution on ibm\_boston, $10{,}000$ shots per row. The
$k=0$ rows carry no oracle and test the accepted law; the $k=1$ rows add one
amplification round. Exact values are in parentheses.}
\label{tab:eval_hardware}
\footnotesize
\setlength{\tabcolsep}{3pt}
\begin{tabular}{lrrrr}
\toprule
Circuit & 2q gates & Acceptance & TVD & $P(\text{opt})$ \\
\midrule
$k=0$, unit demands        &    8 & 0.100 (0.098) & 0.030 & 0.609 (0.632) \\
$k=0$, demands $(1,2,1,2)$ &    8 & 0.138 (0.141) & 0.049 & 0.436 (0.444) \\
$k=0$, six UEs             &    8 & 0.139 (0.149) & 0.028 & 0.431 (0.423) \\
$k=1$, three UEs           &  364 & 0.337 (0.985) & 0.106 & 0.373 (0.436) \\
$k=1$, demands $(1,2,1,2)$ &  787 & 0.125 (0.838) & 0.159 & 0.299 (0.444) \\
$k=1$, six UEs             & 1868 & 0.036 (0.860) & 0.248 & 0.222 (0.423) \\
\bottomrule
\end{tabular}
\end{table}


% ---------------------------------------------------------------------
% [8] Abstract, one word of scope.  The reconstruction claim holds on
%     hardware for the unamplified sampler; leaving it unqualified invites
%     the question the repository answers.
%     Was: "...exact constraint satisfaction, reconstruction of the
%     distribution over feasible assignments, and scalable coordination..."
% ---------------------------------------------------------------------
The results verify exact constraint satisfaction on every simulated and
executed circuit, reconstruction of the distribution over feasible
assignments, and scalable coordination, and establish a practical foundation
for applying a common quantum computational framework to diverse and coupled
network management tasks.


% ---------------------------------------------------------------------
% [9] Sec. I, naming the devices costs four words and answers the first
%     question a reviewer asks.
%     Was: "The proposed quantum algorithm is executed on a real NISQ
%     processor to verify its operation and exact constraint satisfaction."
% ---------------------------------------------------------------------
The proposed quantum algorithm is executed on superconducting and trapped-ion
processors to verify its operation and exact constraint satisfaction.
